\section[唐風]{唐風}

\subsection[蟋蟀]{蟋蟀}

蟋蟀在堂，歲聿其莫。今我不樂，日月其除。無已大康，職思其居。好樂無荒，良士瞿瞿。

蟋蟀在堂，歲聿其逝。今我不樂，日月其邁。無已大康，職思其外。好樂無荒，良士蹶蹶。

蟋蟀在堂，役車其休。今我不樂，日月其慆。無已大康，職思其憂。好樂無荒，良士休休。

\subsection[山有樞]{山有樞}

山有樞，隰有榆。子有衣裳，弗曳弗婁。子有車馬，弗馳弗驅。宛其死矣，他人是愉。

山有栲，隰有杻。子有廷内，弗洒弗埽。子有鐘鼓，弗鼓弗考。宛其死矣，他人是保。

山有漆，隰有栗。子有酒食，何不日鼓瑟？且以喜樂，且以永日。宛其死矣，他人入室。

\subsection[揚之水]{揚之水}

揚之水，白石鑿鑿，素衣朱襮，從子于沃，既見君子，云何不樂。

揚之水，白石皓皓，素衣朱繡，從子于鵠，既見君子，云何其憂。

揚之水，白石粼粼，我聞有命，不敢以告人！

\subsection[椒聊]{椒聊}

椒聊之實，蕃衍盈升。彼其之子，碩大無朋。椒聊且！遠條且！

椒聊之實，蕃衍盈匊。彼其之子，碩大且篤。椒聊且！遠條且！
